\section{Data}
\label{sec:datos}

\subsection{Tasks: O*NET 30.2}

The unit of rating is the task statement of the Occupational Information Network (O*NET), release 30.2 \citep{onet2026database}.
The release contains \cifra{diseno.onet.tareas} task statements for \cifra{diseno.onet.ocupaciones} O*NET-SOC 2019 occupations, each with its occupation title, and two ratings per task from incumbent surveys: importance (scale 1--5) and frequency (seven categories from yearly or less to hourly or more).
We use the mean importance and the mean frequency category of each task as its weight within its occupation, following the task-weighting convention of \citet{acemoglu2011skills}.
Missing importance or frequency values are imputed with the median, as preregistered.
Each task is rated once; a task that appears under several occupations keeps one rating and several weights.

\subsection{Bridge to ISCO-08}

Most countries report employment in the International Standard Classification of Occupations, ISCO-08 \citep{ilo2012isco}, which has 436 unit groups at four digits.
O*NET uses the United States Standard Occupational Classification (SOC).
We bridge the two in three documented steps: O*NET-SOC 2019 to O*NET-SOC 2010 with the official O*NET Center crosswalk; O*NET-SOC 2010 to six-digit SOC 2010 by truncation; and SOC 2010 to ISCO-08 with the crosswalk of \citet{hardy2018educational}, published by the Institute for Structural Research (IBS).
A SOC code that maps to $k$ ISCO groups splits its weight $1/k$.
The result covers \cifra{diseno.puente.isco} of the 436 unit groups with \cifra{diseno.puente.conjuntos} distinct task sets; the three military groups (0110, 0210, 0310) have no O*NET tasks and are excluded, and seven residual groups without tasks of their own inherit the tasks of their three-digit parent and are flagged as imputed.
Four O*NET-SOC 2019 occupations with no 2010 ancestor are assigned by hand.
Appendix~\ref{app:crosswalk} lists every step, the two SOC codes rescued because the IBS file omits them, the seven imputed groups and the four manual assignments; the code is \texttt{datos/construir\_crosswalk.py}.

The bridge replaces the occupation--task matrix used in the first author's earlier work on augmented human capital \citep{espinal2026augmented, espinal2026measuring}, which mapped SOC to ISCO at one digit and so collapsed \cifra{diseno.previo.ocupaciones} Colombian occupations into \cifra{diseno.previo.conjuntos} task sets.
Section~\ref{sec:limites} returns to the consequences.

\subsection{Secondary task bank: ESCO}

O*NET describes work in the United States.
As a check on that origin, we also score the essential skills and competences that the European Skills, Competences, Qualifications and Occupations classification (ESCO, version 1.2) attaches to each of its occupations \citep{esco2026}.
ESCO maps every occupation to an ISCO-08 unit group by construction, so it reaches the 436 groups without a crosswalk, and it is available in 28 languages.
Its units are skills, not tasks, so we treat it as a robustness bank rather than as the primary source.
The ESCO results are reported in Section~\ref{sec:indice} as a sensitivity check and matter most for agriculture, where O*NET has few tasks (Section~\ref{sec:limites}).

\subsection{Employment by occupation: ILOSTAT}

Country weights come from the ILO statistics database, ILOSTAT, series on employment by sex and occupation \citep{ilostat2026}, downloaded through its SDMX interface.
At one digit, \cifra{diseno.ilostat.paises.1d} countries report employment by ISCO-08 major group and \cifra{diseno.ilostat.paises.1d.desde2019} have a reference year of 2019 or later; one of them (Bhutan) reports only armed-forces and unclassified employment in its latest year, so \cifra{diseno.atlas.paises} countries form the atlas.
At two digits, \cifra{diseno.ilostat.paises.2d} countries report sub-major groups and \cifra{diseno.ilostat.paises.2d.desde2019} have a year of 2019 or later.
We use the finest level each country reports, two digits where the sub-major groups are complete (at least thirty groups whose total is within five per cent of the one-digit total) and one digit otherwise, with the pooling rule of Section~\ref{sec:metodo}: \cifra{diseno.atlas.paises.2d} countries enter at two digits and \cifra{diseno.atlas.paises.1d} at one.
Countries whose latest occupational data predate 2019 are excluded, as preregistered.
Table~\ref{tab:cobertura}, generated by \texttt{analisis/tablas.py}, lists coverage by income group.
The same database supplies the informality rate and mean monthly earnings by occupation where they exist, used in the atlas.

\subsection{Income and adoption}

Income groups follow the World Bank classification for fiscal year 2026, and GDP per capita is at purchasing power parity in constant 2021 international dollars \citep{worldbank2026income}.
Observed adoption comes from the Anthropic Economic Index (AEI), which reports usage of Claude by country per working-age adult in its latest release \citep{anthropic2025economic}.
For convergence tests we use the occupation-level scores of \citet{eloundou2024gpts}, \citet{felten2021occupational}, \citet{webb2020impact}, \citet{eisfeldt2023generative} and \citet{gmyrek2025generative}, each carried to ISCO-08 with the same bridge where it is not already in ISCO-08.

\subsection{Household microdata}

For breakdowns by sex, age and informality we use four household surveys that code occupation at four digits in national adaptations of ISCO-08: the Gran Encuesta Integrada de Hogares of Colombia (GEIH) \citep{dane2026geih}, the Pesquisa Nacional por Amostra de Domic\'ilios Cont\'inua of Brazil (PNADC) \citep{ibge2026pnadc}, the Encuesta Nacional de Ocupaci\'on y Empleo of Mexico (ENOE) \citep{inegi2026enoe} and the Encuesta Nacional de Hogares of Peru (ENAHO) \citep{inei2026enaho}.
For Colombia we also use the 2023--2026 monthly series to relate delegability to wage dynamics (Section~\ref{sec:validacion}).
Every input file is listed with its SHA-256 hash and download date in \texttt{datos/MANIFIESTO.jsonl}.
